\documentclass[11pt,a4paper]{article}

% ---------- Encoding & language ----------
\usepackage[T1]{fontenc}
\usepackage[utf8]{inputenc}
\usepackage[english]{babel}

% ---------- Geometry & typography ----------
\usepackage[a4paper,margin=2.5cm]{geometry}
\usepackage{microtype}
\usepackage{setspace}
\setstretch{1.1}

% ---------- Graphics, tables, math ----------
\usepackage{graphicx}
\usepackage{booktabs}
\usepackage{amsmath,amssymb}
\usepackage{float}
\usepackage{listings}
\usepackage{color}

% ---------- Code Listing Style ----------
\definecolor{dkgreen}{rgb}{0,0.6,0}
\definecolor{gray}{rgb}{0.5,0.5,0.5}
\definecolor{mauve}{rgb}{0.58,0,0.82}

\lstset{frame=tb,
  language=Python,
  aboveskip=3mm,
  belowskip=3mm,
  showstringspaces=false,
  columns=flexible,
  basicstyle={\small\ttfamily},
  numbers=none,
  numberstyle=\tiny\color{gray},
  keywordstyle=\color{blue},
  commentstyle=\color{dkgreen},
  stringstyle=\color{mauve},
  breaklines=true,
  breakatwhitespace=true,
  tabsize=3
}

% ---------- Links & refs ----------
\usepackage[hidelinks]{hyperref}
\usepackage[capitalize,nameinlink]{cleveref}

\begin{document}

% =========================================================
%                 FRONT PAGE
% =========================================================
\begin{titlepage}
  \begin{center}
  \thispagestyle{empty}
  \vspace*{2cm}
  {\LARGE \textbf{Lab Report:}}\\[6pt]
  {\Large \textit{Housing Price Estimation using Linear Regression}}\\[18pt]

\vspace{8cm}

% Note: Ensure HHlogo.jpg is in the directory
\includegraphics[width=95px]{HHlogo.jpg}

\vspace{3cm}
  
  {\large
  Kabir
  \\
  Maths Lab
  \\
  \today
  }
  \end{center}  
  \vfill
\end{titlepage}

% =========================================================
%                 MAIN CONTENT
% =========================================================
\pagenumbering{arabic}
\begin{raggedright}

\section{Problem 2: Housing price estimation}

\subsection{Problem Statement}
In this problem, we are tasked with building a model to estimate housing prices based on a set of attributes, such as the number of rooms, crime rate, and pollution levels. Specifically, we want to find a linear relationship that best fits the data:
\[ y = a_0 + a_1x_1 + a_2x_2 + \dots + a_Nx_N \]
Here, $y$ represents the predicted price (in thousands of dollars), $x_i$ are the various features of the house, and $a_i$ are the weights we need to calculate. The goal is to implement this using the Normal Equation method and evaluate how well it works on unseen data.

\subsection{Background}
To solve this, we rely on the \textbf{Method of Least Squares}. The idea is simple: we want to draw a "line" (or hyperplane) through our data points such that the total squared distance between the actual prices and our predicted prices is minimized. This distance is what we call the "error" or "residual."

\textbf{Linear Regression & Bias}:
We model the price $y$ as a weighted sum of features plus a noise term:
\[ y = \mathbf{x}^T \beta + \epsilon \]
A crucial detail here is the bias term, $a_0$. Realistically, a house with "zero" for all features (zero rooms, zero crime, etc.) doesn't have a price of zero—that wouldn't make sense. The regression line needs a starting point, or intercept, which is why we add a column of 1s to our data. This allows the model to learn that baseline price.

\textbf{Derivation of the Normal Equation}:
We start with the cost function $J(\beta)$, which represents the sum of squared errors between our predictions ($X\beta$) and the actual values ($y$):
\[ J(\beta) = ||X\beta - y||^2 = (X\beta - y)^T(X\beta - y) \]

To find the best weights $\beta$, we need to minimize this cost. We do this by expanding the equation and then taking the derivative.

\textit{Step 1: Expand the expression}
\[ J(\beta) = (\beta^T X^T - y^T)(X\beta - y) \]
\[ J(\beta) = \beta^T X^T X \beta - 2\beta^T X^T y + y^T y \]

\textit{Step 2: Take the gradient}
We take the gradient of $J(\beta)$ with respect to $\beta$ and set it to zero to find the minimum:
\[ \nabla_\beta J(\beta) = \frac{\partial}{\partial \beta} (\beta^T X^T X \beta - 2\beta^T X^T y) \]
\[ \nabla_\beta J(\beta) = 2X^T X \beta - 2X^T y \]

\textit{Step 3: Solve for zero}
Setting the gradient to zero:
\[ 2X^T X \beta - 2X^T y = 0 \]
\[ X^T X \beta = X^T y \]

\textit{Step 4: The Solution}
Finally, we multiply by the inverse of $(X^T X)$ to isolate $\beta$:
\[ \beta = (X^T X)^{-1} X^T y \]
This equation gives us the exact mathematical solution for the best-fitting line.

\subsection{Method}
Our approach was straightforward. We loaded the dataset and split it into two parts: 80\% for training the model and 20\% for testing it. This ensures we aren't just memorizing the data but actually learning patterns that generalize.

We repeated this process 100 times, shuffling the data each time. This is important because a single lucky (or unlucky) split could give misleading results. Averaging over 100 runs gives us a much more reliable picture of the model's performance.

\begin{lstlisting}[language=Python, caption=Pseudocode for Housing Price Estimation]
# 1. Load Data
data = load_dataset("Problem2_Dataset.csv")
X = data_features
y = data_target

# 2. Split Data (Training/Testing)
X_train, y_train = subset(X, y, 80%)
X_test, y_test = subset(X, y, 20%)

# 3. Add Bias Term (Intercept)
# Append a column of 1s to X
X_train = augment_with_ones(X_train)
X_test = augment_with_ones(X_test)

# 4. Solve Normal Equation (minimize ||Xb - y||^2)
# Direct formula for beta
beta = inverse(X_train.T @ X_train) @ (X_train.T @ y_train)

# 5. Predict and Evaluate
y_pred = X_test @ beta
rmse = sqrt(mean((y_pred - y_test)^2))
\end{lstlisting}

\subsection{Results}
After running our simulation 100 times, we gathered the following statistics. The target values in this dataset are in \textbf{$\$1000$s}, so the errors should be interpreted in that context.

\textbf{Quantitative Analysis}:
\begin{itemize}
    \item \textbf{Mean RMSE: $\approx 5.04$} \\
    On average, our model's predictions are off by about **$\$5,040$**. Considering the range of housing prices, this is a reasonable starting point, though not perfect.
    
    \item \textbf{RMSE Std Dev: $\approx 0.69$} \\
    The standard deviation is fairly low ($\$690$), which is good news. It means our model is consistent; it doesn't work great one time and terribly the next. It's stable.

    \item \textbf{Best RMSE: $\approx 3.63$} \\
    In the best-case scenario (a particularly "easy" split of the data), the error dropped to around **$\$3,630$**.
\end{itemize}

\textbf{Visual Analysis}:
The graphs below help visualize what's happening.

\begin{figure}[H]
    \centering
    \includegraphics[width=0.7\textwidth]{ne_rmse_histogram.png}
    \caption{RMSE Distribution over 100 runs. You can see the bell curve shape centered around 5.0, confirming the stability we observed in the statistics.}
\end{figure}

\begin{figure}[H]
    \centering
    \includegraphics[width=0.8\textwidth]{ne_coefficients.png}
    \caption{Feature Importance. The bars show which features matter most. \textbf{RM} (rooms) has a large positive bar—more rooms means higher price. \textbf{NOX} (pollution) has a negative bar—cleaner air is worth more. The error bars are small, meaning these relationships are consistent.}
\end{figure}

\begin{figure}[H]
    \centering
    \includegraphics[width=0.7\textwidth]{ne_best_model_fit.png}
    \caption{Best Model Predictions. Ideally, all dots would land on the red dashed line. Our model does well for the most part, but notice how it struggles with the highest-priced homes (top right). This suggests there might be some non-linear factors at play for luxury properties.}
\end{figure}

\subsection{Discussion}
The Normal Equation was an effective tool for this problem. It's exact and easy to implement. However, we should be honest about its limitations in a real-world engineering context.

First of all, it's computationally heavy for large datasets. Inverting a matrix becomes incredibly slow as the number of features grows. If we had 100,000 features instead of 13, this method would crash most computers.

Secondly, it handles outliers poorly. Since we are minimizing the \textit{squared} error, a single massive outlier can pull the regression line disproportionately, skewing the results for everyone else.

\textbf{Future Improvements}:
If we were to take this further, I would suggest:
\begin{itemize}
    \item **Switching to Gradient Descent** if the dataset grows larger. It's an iterative method that scales much better.
    \item **Using Regularization (Ridge/Lasso)**. This would help prevent overfitting and handle situations where features are highly correlated (like TAX and RAD in this dataset).
    \item **Feature Engineering**. We could try combining features or using polynomial terms (e.g., $RM^2$) to capture non-linear relationships, which might fix the improved accuracy on high-end homes.
\end{itemize}

\end{raggedright}

% ---- Appendix ----
\newpage
\section{Appendix}
\subsection{part2.py (Logic)}
\lstinputlisting{part2.py}
\subsection{run\_multiple\_part\_2.py (Runner)}
\lstinputlisting{run_multiple_part_2.py}

\end{document}
